\documentclass[aspectratio=169, 10pt]{beamer}
\usepackage[fontset=fandol]{ctex}
\usepackage{amsmath, amssymb, mathtools}
\usepackage{booktabs}
\usepackage{graphicx}
\usepackage{xcolor}
\usepackage{tikz}
\usetikzlibrary{arrows.meta, positioning, shapes.geometric}

% -- Theme ---------------------------------------------------------------
\usetheme{Madrid}
\usecolortheme{dolphin}
\usefonttheme{professionalfonts}
\setbeamertemplate{navigation symbols}{}

% -- Compact spacing ------------------------------------------------------
\setbeamersize{text margin left=6mm, text margin right=6mm}
\renewcommand{\baselinestretch}{0.95}
\setbeamertemplate{block begin}{\vskip0.8em\begin{beamercolorbox}[ht=2.5ex,dp=0.8ex,leftskip=0.5em,rightskip=0.5em]{block title}\usebeamerfont*{block title}\insertblocktitle\end{beamercolorbox}\vskip0.2ex\begin{beamercolorbox}[dp=0.8ex,leftskip=0.5em,rightskip=0.5em]{block body}\usebeamerfont{block body}}
\setbeamertemplate{block end}{\end{beamercolorbox}\vskip0.8em}
\setbeamertemplate{block alerted begin}{\vskip0.8em\begin{beamercolorbox}[ht=2.5ex,dp=0.8ex,leftskip=0.5em,rightskip=0.5em]{block title alerted}\usebeamerfont*{block title}\insertblocktitle\end{beamercolorbox}\vskip0.2ex\begin{beamercolorbox}[dp=0.8ex,leftskip=0.5em,rightskip=0.5em]{block body alerted}\usebeamerfont{block body}}
\setbeamertemplate{block alerted end}{\end{beamercolorbox}\vskip0.8em}

% -- Colors --------------------------------------------------------------
\definecolor{myblue}{HTML}{2B5EA7}
\definecolor{accentgreen}{HTML}{2E8B57}
\definecolor{accentred}{HTML}{C0392B}
\definecolor{lightbg}{HTML}{F5F7FA}
\setbeamercolor{frametitle}{bg=myblue, fg=white}
\setbeamercolor{structure}{fg=myblue}
\setbeamercolor{block title}{bg=myblue!85, fg=white}
\setbeamercolor{block body}{bg=myblue!8}
\setbeamercolor{alerted text}{fg=accentred}

% -- Auto section title page ---------------------------------------------
\AtBeginSection[]{
  \begin{frame}
    \vfill\centering
    \usebeamerfont{title}\textcolor{myblue}{\insertsectionhead}\par
    \vfill
  \end{frame}
}

% -- Footer with GitHub URL ----------------------------------------------
\setbeamertemplate{footline}{%
  \leavevmode%
  \hbox{%
  \begin{beamercolorbox}[wd=.55\paperwidth,ht=2.5ex,dp=1ex,leftskip=2ex]{author in head/foot}%
    \tiny\texttt{github.com/dechang64/organoid-fl}%
  \end{beamercolorbox}%
  \begin{beamercolorbox}[wd=.35\paperwidth,ht=2.5ex,dp=1ex,center]{author in head/foot}%
    \tiny\insertshorttitle%
  \end{beamercolorbox}%
  \begin{beamercolorbox}[wd=.10\paperwidth,ht=2.5ex,dp=1ex,center]{author in head/foot}%
    \tiny\insertframenumber{} / \inserttotalframenumber%
  \end{beamercolorbox}}%
  \vskip0pt%
}

% -- Metadata ------------------------------------------------------------
\title[Organoid-FL]{Organoid-FL: 类器官检测联邦学习平台}
\subtitle{项目进展汇报}
\date{2026年6月}

\begin{document}

% ========================================================================
\begin{frame}[plain]
  \titlepage
\end{frame}

\begin{frame}{目录}
  \tableofcontents
\end{frame}

% ========================================================================
\section{问题与背景}

\begin{frame}{什么是类器官？}
  \begin{columns}[T]
    \column{0.55\textwidth}
    \begin{block}{类器官（Organoid）}
      \begin{itemize}
        \item 体外培养的微型器官（肺、肠、脑等）
        \item 用于药物筛选、个性化医疗、疾病建模
        \item 显微镜下成像后需\textbf{人工计数和形态分析}
        \item 通量低、主观性强、跨实验室一致性差
      \end{itemize}
    \end{block}
    \vspace{0em}
    \begin{alertblock}{核心痛点}
      类器官图像分析依赖人工，成为药物研发管线的瓶颈。
      自动化检测是解决方案，但面临多个技术挑战。
    \end{alertblock}
    \column{0.40\textwidth}
    \begin{block}{技术挑战}
      \begin{itemize}
        \item \alert{超大图}：6K$\times$5.7K 像素
        \item \alert{密集小目标}：单图数百个 organoid
        \item \alert{标注噪声}：多标注者不一致
        \item \alert{数据分散}：多中心数据无法共享
      \end{itemize}
    \end{block}
  \end{columns}
\end{frame}

\begin{frame}{为什么需要联邦学习？}
  \begin{block}{数据隐私与数据分散的矛盾}
    \begin{itemize}
      \item 多家医院/实验室各自有类器官图像数据
      \item 涉及患者来源 → \alert{无法集中共享}
      \item 单中心数据量不足 → 模型泛化性差
      \item 各中心培养条件/成像设备不同 → 数据分布异构（non-IID）
    \end{itemize}
  \end{block}
  \vspace{0em}
  \begin{block}{联邦学习（Federated Learning）}
    \begin{itemize}
      \item \textbf{数据不动，模型动}：各中心本地训练，只上传模型参数
      \item \textbf{数据不动，知识动}：聚合后的全局模型吸收所有中心的知识
      \item \textbf{数据不动，推理动}：全局模型下发给各中心本地推理
    \end{itemize}
  \end{block}
  \vspace{0em}
  \centering\alert{FL + Organoid 是未被探索的交叉方向}
\end{frame}

% ========================================================================
\section{目标与交付}

\begin{frame}{项目目标}
  \begin{block}{总目标}
    构建一个\textbf{隐私保护的类器官图像分析联邦学习平台}，
    实现多中心协同训练，在检测精度上超越单中心 SOTA。
  \end{block}
  \vspace{0.5em}
  \begin{block}{具体指标}
    \begin{itemize}
      \item[\textbf{R1}] \textbf{检测精度}：MultiOrg 数据集 mAP@0.5 超越 SOTA（SSD 68.1\%）
      \item[\textbf{R2}] \textbf{FL 框架}：支持多中心 non-IID 场景的检测模型联邦训练
      \item[\textbf{R3}] \textbf{聚合策略}：EWA 熵加权聚合，有效区间理论指导
      \item[\textbf{R4}] \textbf{评估方法论}：标注者选择 + 评估方式的系统化框架
      \item[\textbf{R5}] \textbf{论文/专利}：FL+Organoid 方向（未被探索的交叉领域）
    \end{itemize}
  \end{block}
\end{frame}

\begin{frame}{交付物清单}
  \begin{columns}[T]
    \column{0.48\textwidth}
    \begin{block}{代码与平台}
      \begin{itemize}
        \item organoid-fl 开源仓库
        \item 数据预处理管线（tiling + 多标注者）
        \item FL 训练框架（FedAvg + EWA + FedProx）
        \item SAHI 全图推理评估脚本
        \item Streamlit 可视化平台
      \end{itemize}
    \end{block}
    \column{0.48\textwidth}
    \begin{block}{研究产出}
      \begin{itemize}
        \item SOTA 调研报告（4份）
        \item 评估方法论（标注者影响量化）
        \item EWA Effective Interval 理论
        \item 论文草稿（FL+Organoid）
        \item 专利布局（3个方向）
      \end{itemize}
    \end{block}
  \end{columns}
\end{frame}

% ========================================================================
\section{SOTA 调研}

\begin{frame}{类器官检测 SOTA 全景}
  \begin{block}{数据集与 SOTA}
    \begin{tabular}{llcc}
      \toprule
      数据集 & SOTA 方法 & mAP@0.5 & 年份 \\
      \midrule
      Intestinal (840张/4类) & Deliod (YOLOv8s 改) & 87.5\% & 2025 \\
      MultiOrg (411张/1类) & SSD & 68.1\% & 2024 \\
      Colorectal Cancer & YOLOv4 & 36.9\% & 2023 \\
      Gastric Organoid & YOLOv5 & 79\% acc & 2023 \\
      \bottomrule
    \end{tabular}
  \end{block}
  \vspace{0em}
  \begin{block}{关键发现}
    \begin{itemize}
      \item \textbf{小目标检测是核心瓶颈}：所有改进都围绕小目标（+11.6pp vs 分类+0.7pp）
      \item \textbf{数据增强贡献 > 架构改进}：Orga-Dete 数据增强 +3.5pp，BiFPN+MPCA +1.5pp
      \item \textbf{NMS 误删密集目标}：传统 NMS 对密集 organoid 场景有结构性缺陷
      \item \textbf{FL + Organoid 无已发表研究}：无论文、无专利
    \end{itemize}
  \end{block}
\end{frame}

\begin{frame}{SOTA 检测器对比（2025--2026）}
  \begin{block}{通用检测 SOTA}
    \begin{tabular}{llcll}
      \toprule
      模型 & 类型 & COCO mAP & 优势 & 来源 \\
      \midrule
      YOLOv12 & CNN+Attention & 55.9\% & 速度最优 & Tian et al., 2025 \\
      RF-DETR & Transformer & \textbf{60+}\% & NMS-free & Solawetz et al., 2025 \\
      D-FINE & Transformer & 59.5\% & 定位精度 & Peng et al., ICLR 2025 \\
      RT-DETRv3 & Transformer & 55.3\% & 实时DETR & Lv et al., 2025 \\
      \bottomrule
    \end{tabular}
  \end{block}
  \vspace{0em}
  \begin{block}{Organoid 检测 SOTA}
    \begin{tabular}{llll}
      \toprule
      工作 & 数据集 & 模型 & mAP@0.5 \\
      \midrule
      Bukas et al. (NeurIPS 2024) & MultiOrg & SSD & 68.1\% \\
      Yang et al. (Biomimetics 2025) & Mouse Organoid & RTMDet+IDP-Head & +5pp vs baseline \\
      Orga-Dete (Appl. Sci. 2025) & Lung Organoid & YOLOv8-Lite & TIDE 分析 \\
      MEF-YOLO (Frontiers 2025) & Organoid & YOLOv10n+MEF & fine-grained \\
      \textbf{Ours} & MultiOrg & YOLOv12s+SAHI & \textbf{68.8\%} \\
      \bottomrule
    \end{tabular}
  \end{block}
  \vspace{0em}
  \begin{alertblock}{RF-DETR 的结构性优势}
    NMS-free 架构在密集 organoid 场景无 NMS 误删问题；单类检测 94.6\% > YOLOv12X 92.5\%（COCO）
  \end{alertblock}
\end{frame}

\begin{frame}{SAHI 优化的理论基础}
  \begin{block}{FP 来源分析}
    \textbf{Orga-Dete TIDE 分解}（Appl. Sci. 2025）：背景误检是 organoid 检测最大误差源，滑动窗口重叠区域加剧此问题。
  \end{block}
  \vspace{0em}
  \begin{block}{优化策略与文献依据}
    \begin{itemize}
      \item \textbf{Overlap}（YOLO-Patch-Based-Inference, 2025）：建议 15\%--40\%；本实验确认 0.3 最优
      \item \textbf{Soft-NMS}（Bodla et al., ICCV 2017）：高斯衰减替代硬删除，保留更多 TP
      \item \textbf{WBF vs NMS}（Solovyev et al.）：WBF 融合优于 NMS，但 Soft-NMS 在 AP 上更优
      \item \textbf{论文协议}（Bukas et al., NeurIPS 2024）：512+2048 双尺度 + NMS 0.5
      \item \textbf{负样本}（Bulten et al., PMC 2021）：空白 patch，FP 62.5\%$\rightarrow$25\%
    \end{itemize}
  \end{block}
\end{frame}

\begin{frame}{MultiOrg 数据集详解}
  \begin{columns}[T]
    \column{0.50\textwidth}
    \begin{block}{数据特点}
      \begin{itemize}
        \item 411 张肺类器官 2D 显微镜图
        \item 6K$\times$5.7K 像素（超大图）
        \item 63,042 个 bbox 标注
        \item \textbf{单类} organoid（论文设定）
        \item Train 356 / Test 55
      \end{itemize}
    \end{block}
    \column{0.46\textwidth}
    \begin{block}{多标注者结构}
      \begin{itemize}
        \item \textbf{Train}：每图 1 个标注者（A 或 B）
        \item \textbf{Test}：每图 3 套标注
        \begin{itemize}
          \item t0：初始标注（噪声高）
          \item t1\_A：二次标注（精简）
          \item t1\_B：二次标注（最精准）
        \end{itemize}
        \item 标注者差异 → 评估方法论的关键变量
      \end{itemize}
    \end{block}
  \end{columns}
  \vspace{0em}
  \begin{alertblock}{论文 SOTA 基线（Table 3）}
    SSD 68.1\% mAP@0.5 \quad | \quad YOLOv3 64.6\% \quad | \quad RTMDet 61.6\% \quad | \quad Faster R-CNN 60.7\%
  \end{alertblock}
\end{frame}

% ========================================================================
\section{方法论}

\begin{frame}{技术路线}
  \centering
  \begin{tikzpicture}[
    node distance=0.8cm and 1.2cm,
    box/.style={draw, rounded corners, minimum width=2.8cm, minimum height=0.8cm,
                font=\small, align=center, fill=myblue!10},
    arrow/.style={-Stealth, thick, myblue}
  ]
    \node[box] (data) {数据预处理\\Tiling + 多标注者};
    \node[box, right=of data] (detect) {检测训练\\YOLOv12 + freebies};
    \node[box, right=of detect] (eval) {评估\\Patch级 + SAHI全图};
    \node[box, below=of detect] (fl) {联邦学习\\FedAvg + EWA + FedProx};
    \node[box, right=of fl] (deploy) {平台\\Streamlit + Rust HNSW};
    
    \draw[arrow] (data) -- (detect);
    \draw[arrow] (detect) -- (eval);
    \draw[arrow] (detect) -- (fl);
    \draw[arrow] (fl) -- (deploy);
    \draw[arrow] (eval) -- (deploy);
  \end{tikzpicture}
  \vspace{0.5em}
  \begin{itemize}
    \item \textbf{第一阶段}：单中心检测精度超越 SOTA（当前重点）
    \item \textbf{第二阶段}：FL 框架验证（EWA Effective Interval）
    \item \textbf{第三阶段}：FL + Organoid 系统集成
  \end{itemize}
\end{frame}

\begin{frame}{数据预处理：从错误到正确}
  \begin{block}{三个根本性错误的修正}
    \begin{tabular}{p{3.5cm}p{4cm}p{4cm}}
      \toprule
      & \textbf{错误做法} & \textbf{正确做法（论文方法）} \\
      \midrule
      类别 & 2类（Normal/Macros） & \textbf{单类} organoid \\
      边界 bbox & 裁剪到 patch 边缘 & \textbf{丢弃}边界 bbox \\
      评估 & patch 上直接评估 & 全图滑动窗口推理 \\
      \bottomrule
    \end{tabular}
  \end{block}
  \vspace{0em}
  \begin{alertblock}{Tip: 使用新数据集前必须仔细读论文}
    这三个错误都是因为没有仔细阅读 MultiOrg 论文（NeurIPS 2024）导致的。论文明确说明：单类设定、边界 bbox 处理方式、全图滑动窗口评估协议。\textbf{先读文档再动手，不要猜}。
  \end{alertblock}
  \vspace{0em}
  \begin{block}{标注者选择 bug（本次新发现）}
    Test set 有 3 套标注（t0/t1\_A/t1\_B），代码合并后 \alert{GT 重复 3 倍}。修复后 mAP@0.5 从 \textbf{40.6\% → 66.9\%}（+26.3pp）。
  \end{block}
\end{frame}

% ========================================================================
\section{实验结果}

\begin{frame}{标注者选择对评估的影响}
  \begin{block}{同一模型，不同标注者评估}
    \begin{tabular}{lcccc}
      \toprule
      标注者 & Instances & mAP@0.5 & mAP@0.5:0.95 & Recall \\
      \midrule
      t0（初始） & 6465 & 66.9\% & 32.6\% & 60.2\% \\
      t1\_A（二次） & 3950 & 70.5\% & 38.7\% & 64.8\% \\
      \textbf{t1\_B（最精准）} & 4112 & \textbf{78.9\%} & \textbf{53.9\%} & \textbf{72.3\%} \\
      \bottomrule
    \end{tabular}
  \end{block}
  \vspace{0em}
  \begin{alertblock}{关键发现}
    \begin{itemize}
      \item 标注者选择可造成 \textbf{12pp mAP@0.5 差异}（66.9\% vs 78.9\%）
      \item t0 噪声最大（6465 instances，含误标）；t1\_B 最精简（4112 instances）
      \item \textbf{评估方法论和模型本身一样重要}——噪声标注会严重低估模型能力
    \end{itemize}
  \end{alertblock}
\end{frame}

\begin{frame}{v5 训练结果}
  \begin{columns}[T]
    \column{0.55\textwidth}
    \begin{block}{配置}
      \begin{itemize}
        \item 模型：YOLOv12s（9.2M 参数）
        \item 分辨率：640$\times$640
        \item Freebies：copy\_paste=1.0, mixup=0.15, \\
              label\_smoothing=0.1, cos\_lr, close\_mosaic=15
        \item Epochs：200（74 epoch 停止）
        \item 硬件：RTX 3060 12GB
      \end{itemize}
    \end{block}
    \column{0.40\textwidth}
    \begin{block}{训练趋势（t1\_B 评估）}
      \begin{tabular}{cc}
        \toprule
        Epoch & mAP@0.5 \\
        \midrule
        1 & 73.7\% \\
        10 & \textbf{78.1\%}（best） \\
        55 & 77.8\% \\
        74 & 77.5\%（停止） \\
        \bottomrule
      \end{tabular}
      \vspace{0em}\\
      \tiny ep10 达峰后平台震荡
    \end{block}
  \end{columns}
  \vspace{0em}
  \begin{block}{对比 v4（无 freebies，66 epoch）}
    v4 最终 mAP@0.5 = 66.9\%（t0）\quad→\quad v5 ep10 = 78.1\%（t1\_B），\textbf{+11.2pp}
  \end{block}
\end{frame}

\begin{frame}{SAHI 全图推理优化}
  \begin{block}{问题：Patch 级 vs 全图推理差距}
    Patch 级 78.1\% → SAHI 全图 64.8\%（$-$13.3pp），核心瓶颈是 \alert{FP 暴增}（4463 FP vs 4114 TP）
  \end{block}
  \vspace{0em}
  \begin{block}{参数优化实验（同一 best.pt）}
    \begin{tabular}{lcccc}
      \toprule
      配置 & mAP@0.5 & Precision & Recall & FP \\
      \midrule
      WBF, overlap=0.5 & 64.8\% & 48.0\% & 83.0\% & 4463 \\
      NMS, overlap=0.5 & 62.4\% & 42.8\% & 86.3\% & 5720 \\
      WBF, overlap=0.3 & 66.1\% & 52.8\% & 81.3\% & 3602 \\
      \textbf{Soft-NMS, overlap=0.3} & \textbf{68.8\%} & 27.0\% & 85.5\% & 11466 \\
      \midrule
      Soft-NMS + score filter 0.3 & 66.5\% & \textbf{52.1\%} & 81.1\% & 3690 \\
      \bottomrule
    \end{tabular}
  \end{block}
  \vspace{0em}
  \begin{alertblock}{最优配置超 SOTA}
    \begin{itemize}
      \item \textbf{论文报告}：Soft-NMS + overlap=0.3 → mAP@0.5 = \textbf{68.8\%}，超 SSD SOTA \textbf{+0.7pp}
      \item \textbf{实际部署}：+ score filter 0.3 → Precision 52.1\%，Recall 81.1\%（可用平衡）
    \end{itemize}
  \end{alertblock}
\end{frame}

\begin{frame}{SAHI 优化的理论基础}
  \begin{block}{FP 来源分析}
    \begin{itemize}
      \item \textbf{Orga-Dete TIDE 分解}（Applied Sciences 2025）：organoid 检测中背景误检（Bkg Error）是最大误差源，滑动窗口重叠区域加剧此问题
      \item \textbf{碎片框}：overlap 区域目标被多个窗口各检测一次，NMS 删除不彻底
    \end{itemize}
  \end{block}
  \vspace{0em}
  \begin{block}{优化策略与文献依据}
    \begin{itemize}
      \item \textbf{Overlap 优化}（YOLO-Patch-Based-Inference, 2025）：建议 15\%--40\%，过高导致重复检测；本实验确认 0.3 最优
      \item \textbf{Soft-NMS}（Bodla et al., ICCV 2017）：高斯衰减替代硬删除，密集场景保留更多 TP
      \item \textbf{NMS vs WBF}（Solovyev et al.）：WBF 加权融合优于 NMS 硬删除，但 Soft-NMS 在 AP 指标上更优
      \item \textbf{MultiOrg 论文协议}（Bukas et al., NeurIPS 2024）：512+2048 双尺度 + downsample 2/8 + NMS 0.5
    \end{itemize}
  \end{block}
\end{frame}

\begin{frame}{当前 vs SOTA 定位}
  \begin{block}{MultiOrg 数据集 mAP@0.5 对比}
    \begin{tabular}{lcc}
      \toprule
      配置 & mAP@0.5 & vs SOTA \\
      \midrule
      Faster R-CNN（论文） & 60.7\% & $-$7.4pp \\
      RTMDet（论文） & 61.6\% & $-$6.5pp \\
      YOLOv3（论文） & 64.6\% & $-$3.5pp \\
      \textbf{SSD（SOTA）} & \textbf{68.1\%} & — \\
      \midrule
      v4 patch 级（t0） & 66.9\% & $-$1.2pp \\
      v4 patch 级（t1\_B） & 78.9\% & +10.8pp \\
      v5 patch 级（t1\_B） & 78.1\% & +10.0pp \\
      \textbf{v5 SAHI 全图（t1\_B）} & \textbf{68.8\%} & \textbf{+0.7pp} \\
      \bottomrule
    \end{tabular}
  \end{block}
  \vspace{0em}
  \begin{alertblock}{结论}
    \begin{itemize}
      \item Patch 级超 SOTA \textbf{+10.0pp}（78.1\% vs 68.1\%）
      \item SAHI 全图推理超 SOTA \textbf{+0.7pp}（68.8\% vs 68.1\%）
      \item Soft-NMS + overlap=0.3 是当前最优推理配置
    \end{itemize}
  \end{alertblock}
\end{frame}

% ========================================================================
\section{联邦学习}

\begin{frame}{为什么要做 FL + Organoid？}
  \begin{block}{现实需求}
    \begin{itemize}
      \item 多中心类器官数据天然分散（各医院/实验室独立培养）
      \item 培养条件、成像设备、标注标准不同 → 数据分布异构（non-IID）
      \item 患者来源数据涉及隐私 → 无法集中共享
      \item 单中心数据量不足 → 模型泛化性差
    \end{itemize}
  \end{block}
  \vspace{0em}
  \begin{block}{未被探索的交叉领域}
    \begin{itemize}
      \item \textbf{论文}：FL + Organoid 无任何已发表论文
      \item \textbf{专利}：FL + 类器官检测系统无专利
      \item \textbf{技术挑战}：检测模型的 FL（vs 分类 FL）研究尚少
    \end{itemize}
  \end{block}
  \vspace{0em}
  \centering\alert{FL 是刚需（数据天然分散），不是可选附加}
\end{frame}

\begin{frame}{EWA 熵加权聚合策略}
  \begin{block}{核心思想}
    传统 FedAvg 按样本量均匀加权，忽略客户端模型质量差异。\\
    EWA 用验证集熵作为信号，\textbf{给更好的客户端更高权重}：
    \[
      w_c = \frac{\exp(-H_c / \alpha)}{\sum_{c'} \exp(-H_{c'} / \alpha)}
    \]
    其中 $H_c$ 是客户端 $c$ 在验证集上的预测熵，$\alpha$ 是温度参数。
  \end{block}
  \vspace{0em}
  \begin{block}{EWA Effective Interval 理论}
    \begin{itemize}
      \item \textbf{IID 场景}：客户端质量差异小 → EWA 退化为均匀加权 → \alert{无优势}
      \item \textbf{适度异质}：质量差异可区分 → EWA 优势最明显 → \alert{+1.4\textasciitilde 1.5pp}
      \item \textbf{极端异质}：信号噪声大 → EWA 略优于 FedAvg → \alert{+1.1pp}
      \item 实验验证：3-client non-IID，EWA 在适度异质下最佳
    \end{itemize}
  \end{block}
\end{frame}

% ========================================================================
\section{未来计划}

\begin{frame}{下一步：检测精度优化}
  \begin{block}{近期（一至两周）}
    \begin{itemize}
      \item \textbf{RF-DETR 对比实验}：NMS-free Transformer 架构，单类检测 94.6\% > YOLOv12X 92.5\%（COCO）
      \item Soft-NMS 后 score 过滤提升 Precision（当前 27\%）
      \item 负样本训练：加入空白 patch 减少 FP（Bulten et al., PMC 2021：FP 62.5\%→25\%）
    \end{itemize}
  \end{block}
  \vspace{0em}
  \begin{block}{中期（一至二月）}
    \begin{itemize}
      \item \textbf{多标注者标签融合}：t1\_A + t1\_B 共识标签（IoU 投票）
      \item \textbf{标签噪声处理}：CLOD 清洗 + ELD 蒸馏
      \item \textbf{FL 实验}：study-based non-IID，3-client EWA vs FedAvg vs FedProx
    \end{itemize}
  \end{block}
\end{frame}

\begin{frame}{下一步：论文与专利}
  \begin{columns}[T]
    \column{0.48\textwidth}
    \begin{block}{论文方向}
      \begin{itemize}
        \item \textbf{FL + Organoid 检测系统}\\
              首个该方向的系统性研究
        \item \textbf{EWA Effective Interval}\\
              FL 聚合策略的有效条件理论
        \item \textbf{评估方法论}\\
              标注者选择对检测评估的影响
      \end{itemize}
    \end{block}
    \column{0.48\textwidth}
    \begin{block}{专利方向}
      \begin{itemize}
        \item \textbf{FL + Organoid 检测系统}\\
              无已发表研究
        \item \textbf{多标注者共识 + organoid}\\
              标签清洗方法
        \item \textbf{RF-DETR + Organoid 推理}\\
              NMS-free 密集场景
      \end{itemize}
    \end{block}
  \end{columns}
\end{frame}

% ========================================================================
\section{未来方向：视觉原语框架}

\begin{frame}{从检测优化到范式转换}
  \begin{block}{当前范式的根本局限}
    端到端"像素$\rightarrow$bbox"把概念固化在权重中。换领域/标注者/中心 $\rightarrow$ 权重失效。SAHI/NMS/Soft-NMS 都是后处理治标。
  \end{block}
  \vspace{0em}
  \begin{alertblock}{启发：Thinking with Visual Primitives (DeepSeek-AI, 2025)}
    将视觉原语（点、bbox）作为推理最小单元，解决 Reference Gap（语言指代空间有歧义），保留指代不确定性。
  \end{alertblock}
  \vspace{0em}
  \begin{block}{映射到我们的问题}
    \begin{itemize}
      \item Soft-NMS 赢 6pp $\leftarrow$ 保留"指代不确定性"
      \item 多标注者 = 多提案，不需选"正确答案"
      \item FL 泛化 = 聚合视觉原语分布，不是聚合权重
    \end{itemize}
  \end{block}
\end{frame}

\begin{frame}{三层通用视觉框架设想}
  \begin{block}{Layer 1: Perception（感知层）}
    SAM2/SegAny 输出无类别 visual primitives（mask+位置+形态）。不学习"这是什么"，只学习"这里有什么"。\textbf{跨域通用}。
  \end{block}
  \vspace{0em}
  \begin{block}{Layer 2: Reasoning（推理层）}
    VLM/LLM 看着 primitives + 全图上下文推理。泛化来源：语义理解，不是像素匹配。推理可解释（CoT）。
  \end{block}
  \vspace{0em}
  \begin{block}{Layer 3: Collaboration（协作层）}
    FL 不聚合权重，聚合 primitive 分布。各中心上传"形态统计"，不上传像素/权重。\textbf{数据不动，知识动}。
  \end{block}
\end{frame}

\begin{frame}{已有探索与技术栈}
  \begin{block}{EWA-Fed: 视觉原语编解码（已实现）}
    \texttt{primitives.py}（BOX/POINT/PATH + Shannon 熵）+ \texttt{aggregator.py}（熵加权聚合）+ \texttt{conformity.py}（从众检测）。当前仅用于监控，未用于推理。
  \end{block}
  \vspace{0em}
  \begin{block}{技术栈映射}
    \begin{tabular}{ll}
      \toprule
      框架层 & 已有实现 \\
      \midrule
      Perception & SAM2 (PCB-Defect-FL), DINOv2 (embodied-fl) \\
      Reasoning & LLM (Claw for Med GPS), RAG (rag\_engine) \\
      Collaboration & FedCtx HNSW + 审计链, EWA 聚合 \\
      \bottomrule
    \end{tabular}
  \end{block}
  \vspace{0em}
  \begin{alertblock}{下一步}
    EWA 从"监控框架"升级为"训练框架"：primitive 传输 + 推理替代检测
  \end{alertblock}
\end{frame}

% ========================================================================
\section{总结}

\begin{frame}{总结}
  \begin{columns}[T]
    \column{0.48\textwidth}
    \begin{block}{已完成}
      \begin{itemize}
        \item[\checkmark] SOTA 全面调研
        \item[\checkmark] 数据预处理管线（tiling + 多标注者）
        \item[\checkmark] 标注 bug 修复（mAP +26pp）
        \item[\checkmark] Patch 级超 SOTA（78.1\%, +10pp）
        \item[\checkmark] SAHI 全图超 SOTA（68.8\%）
        \item[\checkmark] EWA Effective Interval 理论
      \end{itemize}
    \end{block}
    \column{0.48\textwidth}
    \begin{block}{进行中 / 待完成}
      \begin{itemize}
        \item[$\circ$] RF-DETR 对比实验
        \item[$\circ$] 负样本训练
        \item[$\circ$] FL + Organoid 系统实验
        \item[$\circ$] 论文撰写 + 专利申请
      \end{itemize}
    \end{block}
  \end{columns}
\end{frame}

\begin{frame}[allowframebreaks]{参考文献}
  \footnotesize
  \begin{thebibliography}{99}
    \bibitem{bukas2024} Bukas et al. \textit{MultiOrg: A Multi-rater Organoid-detection Dataset}. NeurIPS, 2024.
    \bibitem{tian2025} Tian et al. \textit{YOLOv12: Attention-Centric Real-Time Object Detectors}. arXiv:2502.12524, 2025.
    \bibitem{solawetz2025} Solawetz et al. \textit{RF-DETR: Real-Time DEtection TRansformer}. ICLR, 2026. GitHub: roboflow/rf-detr.
    \bibitem{peng2025} Peng et al. \textit{D-FINE: Redefine Regression Task of DETRs as Fine-grained Distribution Refinement}. ICLR, 2025.
    \bibitem{lv2025} Lv et al. \textit{RT-DETRv3: Revisiting Real-Time Detection Transformer}. arXiv:2602.21010, 2025.
    \bibitem{bodla2017} Bodla et al. \textit{Soft-NMS -- Improving Object Detection with One Line of Code}. ICCV, 2017.
    \bibitem{solovyev2019} Solovyev et al. \textit{Weighted Boxes Fusion: Ensembling Boxes from Different Object Detection Models}. Image and Vision Computing, 2019.
    \bibitem{akyon2022} Akyon et al. \textit{Slicing Aided Hyper Inference (SAHI)}. CVPRW, 2022.
    \bibitem{yang2025idp} Yang et al. \textit{IDP-Head: An Interactive Dual-Perception Architecture for Organoid Detection}. Biomimetics, 2025.
    \bibitem{orgadete2025} Orga-Dete. \textit{An Improved Lightweight Deep Learning Model for Lung Organoid Detection}. Applied Sciences, 2025.
    \bibitem{mefyolo2025} MEF-YOLO. \textit{A Deep Lightweight Model for Solving the Fine-grained of Organoid}. Frontiers in Molecular Biosciences, 2025.
    \bibitem{bulten2021} Bulten et al. \textit{Overcoming the Limitations of Patch-based Learning to Detect Cancer in Whole Slide Images}. PMC, 2021.
    \bibitem{lu2025} Lu et al. \textit{Thinking with Visual Primitives}. DeepSeek-AI, 2025.
    \bibitem{kirillov2023} Kirillov et al. \textit{Segment Anything (SAM)}. ICCV, 2023.
    \bibitem{oquab2024} Oquab et al. \textit{DINOv2: Learning Robust Visual Features without Supervision}. TMLR, 2024.
  \end{thebibliography}
\end{frame}

\begin{frame}[plain]
  \centering
  \vfill
  \Large\textcolor{myblue}{Thank You!}\\[0.8em]
  \normalsize Questions \& Discussion\\[1.5em]
  \footnotesize\texttt{github.com/dechang64/organoid-fl}
  \vfill
\end{frame}
\end{document}
